\renewcommand{\Anno}{2013}

\renewcommand{\NumeroProblema}{2}

\renewcommand{\Testo}{ 
In un gruppo di 9 studenti, tra cui ci sono i tre amici Alberto,
Barbara e Cristina, bisogna sorteggiare 4 fortunati che parteciperanno
ad un soggiorno all'estero.  Per far questo l'organizzatore prepara
un'urna con 9 palline, ciascuna delle quali contiene un numero intero da 1
a 9, non visibile dall'esterno (ovviamente ogni numero sta in una ed
una sola pallina).  Ogni studente pesca una pallina (dunque studenti
diversi pescano palline diverse) ed i 4 studenti con i numeri pi\`{u}
alti saranno i 4 prescelti.

\begin{enumerate}
	\item[(a)]  Determinare la probabilit\`{a} che nessuno dei 3 amici 
	(Alberto, Barbara e Cristina) sia tra i prescelti. 

	\item[(b)] Determinare la probabilit\`{a} che esattamente uno tra
	Alberto, Barbara e Cristina sia tra i prescelti.

	\item[(c)] Dopo che gli studenti hanno estratto le palline, ma non
	ancora visto i numeri in esse contenuti, l'organizzatore del
	soggiorno si fa dare le palline e, senza comunicare i numeri
	estratti, afferma che ``il numero di Alberto \`{e} pi\`{u} grande
	del numero di Barbara, il quale a sua volta \`{e} pi\`{u} grande
	del numero di Cristina''.
	
	Avuta questa informazione, determinare la
	probabilit\`{a} di Alberto, la probabilit\`{a} di Barbara e 
	probabilit\`{a} di Cristina di essere tra i prescelti.
\end{enumerate}
}

\renewcommand{\Risposta}{
\begin{enumerate}
	\renewcommand{\labelenumi}{(\alph{enumi})}
	\item  $5/42$

	\item  $10/21$
	
	\item  $\mbox{Alberto}=37/42$, $\mbox{Barbara}=17/42$, 
	$\mbox{Cristina}=1/21$.
\end{enumerate}
}

\renewcommand{\Aiutino}{
I casi possibili sono sempre tanti quanti i modi di scegliere 3
oggetti su 9.  Perch\'{e}?  
Per quanto riguarda i casi favorevoli, dipende dai vari punti.
\begin{enumerate}
	\renewcommand{\labelenumi}{(\alph{enumi})}
	\item  I casi favorevoli sono i modi di scegliere 3 oggetti su 5. 
	Perch\'{e}?

	\item I casi favorevoli sono i modi di scegliere 1 oggetto su 4 e
	2 oggetti su 5.  Perch\'{e}?
	
	\item Il modo di procedere \`{e} diverso a seconda del personaggio
	che si considera.
	\begin{itemize}
		\item  Nel caso di Alberto, qual \`{e} l'evento complementare? 
	
		\item Nel caso di Cristina, \`{e} come chiedere che i 3 amici
		siano tutti prescelti.  Chi sono i casi favorevoli in questo
		caso?  
	
		\item  Nel caso di Barbara, conviene distinguere il caso in 
		cui tutti sono prescelti ed il caso in cui solo Alberto e 
		Barbara sono prescelti.
	\end{itemize}
\end{enumerate}
}

\renewcommand{\SoluzioneUfficiale}{\TitoloSoluzione 
In tutto l'esercizio possiamo pensare che i casi possibili, tutti 
equiprobabili, siano i modi di scegliere 3 oggetti su 9, cio\`{e} i 3 
numeri sorteggiati dai 3 amici, indipendentemente dall'ordine. I casi 
possibili sono dunque
$$\binom{9}{3}=84.$$

\DomandaA
Se nessuno dei 3 amici viene prescelto, vuol dire che i 3 numeri che 
hanno sorteggiato stanno tra i 5 pi\`{u} piccoli. I casi favorevoli 
sono dunque
$$\binom{5}{3}=10.$$

La probabilit\`{a} richiesta \`{e} quindi $10/84=5/42$.

\DomandaB
Se esattamente uno dei 3 viene prescelto, vuol dire che i 3 numeri 
sorteggiati dai 3 amici sono 1 tra i 4 vincenti, e 2 tra i 5 pi\`{u} 
bassi. I casi favorevoli sono dunque
$$4\cdot\binom{5}{2}=4\cdot 10=40,$$
da cui la probabilit\`{a} richiesta risulta $40/84=10/21$.

\DomandaC
Nel caso di Alberto, consideriamo l'evento complementare, cio\`{e} 
quello in cui Alberto non viene selezionato. Essendo lui il meglio 
messo tra i 3 amici, questo \`{e} equivalente a dire che nessuno dei 
3 viene selezionato. Ora dal punto~(a) sappiamo che ci\`{o} accade 
con probabilit\`{a} 5/42. Passando al complementare, la 
probabilit\`{a} che Alberto sia selezionato sar\`{a} dunque
$$\mbox{Probabilit\`{a}}(\mbox{Alberto})=1-\frac{5}{42}=\frac{37}{42}.$$

Passiamo ora a Cristina, la quale sar\`{a} selezionata se e solo se i 
numeri estratti dai 3 amici stanno tra i 4 pi\`{u} grandi. I casi 
favorevoli sono dunque 
$$\binom{4}{3}=4,$$ 
da cui
$$\mbox{Probabilit\`{a}}(\mbox{Cristina})=\frac{4}{84}=\frac{1}{21}.$$

Consideriamo infine Barbara.  Possiamo dividere i casi in cui viene
selezionata in due sottoclassi: i casi in cui anche Cristina viene
selezionata, che abbiamo gi\`{a} contato e sono 4, e quelli in cui
solo Alberto e Barbara sono selezionati.  Questi ultimi sono tanti
quanti i modi di scegliere 2 oggetti su 4 (i punteggi vincenti di
Albetto e Barbara) ed un oggetto su 5 (il punteggio perdente di
Cristina), e cio\`{e}
$$\binom{4}{2}\cdot 5=6\cdot 5=30.$$

In conclusione avremo che
$$\mbox{Probabilit\`{a}}(\mbox{Barbara})=\frac{4+30}{84}=
\frac{34}{84}=\frac{17}{42}.$$

}

\renewcommand{\SoluzioniAlternative}{
}

\renewcommand{\Osservazioni}{
Nei punti~(a) e~(b) possiamo pensare che i casi possibili, sempre 
equiprobabili, siano i modi di assegnare 9 numeri a 9 persone, dunque 
$9!$. A questo punto, per\`{o}, bisogner\`{a} tenere conto 
dell'ordine anche nel calcolo dei casi favorevoli. Ad esempio, per 
contare i casi favorevoli per il punto~(a), bisogner\`{a} scegliere i 
3 numeri su 5 che spettano ai 3 amici e poi moltiplicare per $3!$ (i 
modi di assegnare i 3 numeri scelti ai 3 amici) e per $6!$ (i modi di 
assegnare i restanti 6 numeri ai restanti 6 studenti).

In poche parole, si finisce per moltiplicare numeratore e 
denominatore per $3!\cdot 6!$, il che ovviamente non cambia i 
risultati.

Nei calcoli del punto~(c), volendo tenere conto dell'ordine, bisogna
moltiplicare numeratori e denominatori per $6!$, perch\'{e} ora
l'ordinamento dei punteggi sui 3 amici \`{e} fisso e si pu\`{o}
scegliere liberamente l'ordinamento dei punteggi solo sui restanti 6
studenti.  }

\renewcommand{\Commenti}{
}